\chapter{Kết luận và hướng phát triển}
\section{Kết quả đạt được}
Project01 đã hiện thực một Diabetes Agent chạy local với mã nguồn độc lập, bốn notebook có output thực, dữ liệu có provenance, model đã huấn luyện và giao diện quan sát được quá trình phối hợp.

Kết quả supervised cho thấy các mô hình học được có khả năng phân biệt nhãn trên test nội bộ cao hơn Dummy. Random Forest được chọn mặc định theo AP validation; MLP vẫn được cung cấp và đánh giá đầy đủ. Truy xuất tìm đúng các ID đích trong bộ câu tự xây, còn bộ hợp đồng và pytest xác nhận những nhánh chức năng đã định. Q-learning học được policy có reward hơn never/rule nhưng chưa vượt fixed trong simulator hiện tại.

Giá trị của thiết kế đa tác tử ở project là trách nhiệm được tách rõ, có thông điệp, kiểm tra và quyền từ chối trước khi tổng hợp. Chưa có kết quả chứng minh kiến trúc này tốt hơn một pipeline đơn dùng cùng công cụ. Việc trình bày giới hạn là cần thiết để tách khả năng chạy được của phần mềm khỏi độ tin cậy trong ứng dụng thực.

\section{Giới hạn hiện tại}
Dữ liệu là khảo sát Hoa Kỳ, nhãn gộp và chưa có thẩm định ngoài mẫu. Mã nhóm nhân khẩu được đưa vào dưới dạng số; chưa kiểm tra thiên lệch theo nhóm hoặc phương án mã hóa khác. Supervised mới có một split/seed và chưa có khoảng tin cậy hoặc external validation.

KB nhỏ, bộ câu hỏi tự xây; Council và luật tiếp nhận chỉ kiểm tra các điều kiện đã định. RL dùng reward tự đặt và chưa có dữ liệu hành vi thật. Demo chưa có các yêu cầu vận hành cho hồ sơ bệnh nhân và không thay thế kế hoạch chăm sóc.

\section{Hướng phát triển có thể kiểm chứng}
Bước tiếp theo nên ưu tiên một bộ câu hỏi độc lập và đánh giá thủ công độ liên quan của trích đoạn, trước khi tăng số agent hoặc bổ sung LLM. Với nhánh dự đoán, cần so sánh mã hóa biến, lặp split/seed, báo khoảng bất định, đánh giá calibration và subgroup trong điều kiện sử dụng xác định.

Với RL, cần kiểm tra độ nhạy theo hệ số reward và động lực phản ứng, đối chiếu với nghiệm tối ưu của MDP nhỏ nếu muốn phân tích sai lệch policy. Chỉ chuyển sang nghiên cứu nhắc việc thực tế khi có dữ liệu, đồng thuận và thiết kế đánh giá phù hợp. Những mở rộng này tạo câu hỏi kiểm chứng được, thay vì chỉ tăng độ phức tạp của hệ thống.
